% Generated draft: requires mathematical and source review.
\documentclass[11pt]{article}
\usepackage[margin=25mm]{geometry}
\usepackage{fontspec}
\setmainfont{Latin Modern Roman}
\usepackage{amsmath,amssymb,mathtools,unicode-math}
\setmathfont{Latin Modern Math}
\usepackage{enumitem,needspace,xurl,hyperref}
\setlength{\parindent}{0pt}
\setlength{\parskip}{5pt}
\setlength{\emergencystretch}{3em}
\newcommand{\N}{\mathbb{N}}
\newcommand{\Z}{\mathbb{Z}}
\newcommand{\Q}{\mathbb{Q}}
\newcommand{\R}{\mathbb{R}}
\newcommand{\C}{\mathbb{C}}
\newcommand{\F}{\mathbb{F}}
\newcommand{\A}{\mathbb{A}}
\newcommand{\PP}{\mathbb{P}}
\newcommand{\E}{\mathbb{E}}
\newcommand{\eps}{\varepsilon}
\newcommand{\dd}{\,\mathrm{d}}
\begin{document}
\section*{UnsolvedMath ID: 800621}
\textbf{Source problem code:} catalogue entry
\subsection{Definitions and mathematical statement}
For purposes of interpreting the displayed formulas, let \(\N^{\N}\) denote the set of functions from \(\N\) to \(\N\). The supplied record instead writes \(Y\subseteq\N\) and \(g\in\N\); the function-space interpretation used here is an inferred correction that requires source verification.

For \(f,g\in\N^{\N}\), define eventual domination by
\[
 f\leq^{*}g\quad\Longleftrightarrow\quad \{n\in\N:f(n)>g(n)\}\text{ is finite}.
\]
A family \(Y\subseteq\N^{\N}\) is dominating if
\[
 (\forall u\in\N^{\N})(\exists f\in Y)\;u\leq^{*}f.
\]
For \(A\subseteq\N^{\N}\), say that \(g\in\N^{\N}\) \emph{avoids middles in \(A\)} if both of the following hold:
\begin{enumerate}
\item For every \(f\in A\), \(g\not\leq^{*}f\).
\item For every \(f,h\in A\), at least one of
\[
 \{n\in\N:f(n)<g(n)\leq h(n)\},\qquad
 \{n\in\N:h(n)<g(n)\leq f(n)\}
\]
is finite.
\end{enumerate}
Using the conventional interpretation of the other invariant, let \(\mathcal M\) be the ideal of meager subsets of \(\R\), where a meager set is a countable union of nowhere dense sets, and put
\[
 \operatorname{cov}(\mathcal M)
 =\min\bigl\{|\mathcal A|:\mathcal A\subseteq\mathcal M,
                  \ \bigcup\mathcal A=\R\bigr\}.
\]
The record does not itself specify this convention.

The supplied question asks whether
\[
 \operatorname{cov}(\mathcal M)
 \leq \operatorname{add}(\mathfrak X,\mathfrak D).
\]
However, its definition of the right-hand side is unresolved. It describes it as the least cardinal \(\kappa\) for which there is a dominating family \(Y\) of cardinality \(\kappa\) such that every partition of \(Y\) into \(\kappa\) pieces has a piece in which no function avoids middles. Under the function-space interpretation above, this requirement cannot hold: the partition into singleton sets is permitted, and for each singleton \(\{f\}\), the function \(g(n)=f(n)+1\) avoids middles in \(\{f\}\). Thus the supplied wording does not define the claimed cardinal. The intended definition of \(\operatorname{add}(\mathfrak X,\mathfrak D)\) must be recovered from the cited source before the displayed question can be treated as a precise mathematical assertion; no replacement definition is asserted here.
\subsection{Short English statement}
The supplied problem asks whether the covering number of the meager ideal is at most a cardinal defined using dominating families and functions that avoid middles. The function domains appear to have lost an exponent, and the stated partition requirement conflicts with the partition into singletons. The intended cardinal therefore requires source clarification. The registry labels the problem open; that label is supplied metadata, not an independently verified status assessment.
\subsection{Sources}
\begin{itemize}
\item Supplied local UnsolvedMath registry record, permanent ID \(800621\), exact source problem code \texttt{catalogue entry}; title: \emph{Selection Principles in Mathematics --- Problem 50 ([ShTb768])}. The record attributes the problem to Tsaban, with proposed year 2006, and labels it open.
\item Tsaban, \emph{Selection Principles and Special Sets of Reals: Open Problems} (2006), Problem 50. Supplied source URL: \url{https://arxiv.org/abs/math/0609601}. The supplied provenance records access on 2026-10-16 and extraction by \texttt{arxiv\_latex}; these provenance claims have not been independently verified here.
\item The supplied background gives attribution/context \texttt{[ShTb768]}, and the statement says that the cardinal is studied there. No bibliographic expansion or separate URL for this reference is supplied, so that attribution remains unresolved.
\end{itemize}
\end{document}
